\chapter{Text: From Bag of Words to Embeddings}\label{ch:ml:text-from-bag-of-words-to-embeddings}

A headline says a company beat estimates, and its shares fall. A dictionary scores the headline positive; a model
trained on returns learns that ``beat estimates but cut guidance'' is bad news; and neither knows which of three companies
called Apex the headline is about. Book~7 (chapter~12) introduced text as data and its first instrument, the dictionary;
Book~8 (chapter~17) traded machine-readable news. This chapter builds the machinery between the words and the forecast:
how a document becomes a vector, how the vector is scored against returns, what topics and embeddings can and cannot
see, and how a mention becomes a security. On a synthetic corpus with a known truth, a return-supervised model reaches
an information coefficient of 0.27 where a dictionary reaches 0.22, embeddings learned without returns reach 0.09 because
they put ``beat'' next to ``missed'', and an entity linker that uses today's tickers maps one mention in eight to the
wrong company or to none.

\section{From documents to vectors}

\begin{definition}[Tokenisation, bag of words, n-gram, tf-idf]\label{def:ml:text-from-bag-of-words-to-embeddings:bow}
\emph{Tokenisation}\index{tokenisation} splits a text into units (words, word pieces, punctuation), usually after
normalising case. A \emph{bag of words}\index{bag of words} represents a document by the counts of its tokens, forgetting
their order. An \emph{n-gram}\index{n-gram} is a sequence of $n$ consecutive tokens; bags of n-grams keep some order.
\emph{Term frequency--inverse document frequency}\index{term frequency--inverse document frequency} (tf-idf) weights the count
of term $w$ in document $d$ by $\log(N/n_w)$, where $n_w$ of the $N$ documents contain $w$, so that words found everywhere
count for little; the rows are then scaled to unit length.
\end{definition}

The chapter's corpus has 20\,000 headlines over six years (1\,512 trading days) about 32 companies: 28 at the start,
three of them called Apex (a software firm, a miner and a healthcare firm), six ticker changes, three renames, and four
delistings whose tickers are reused by new listings 30 days later. Each headline names a company (by its full name, its
short name or its ticker), states one of 19 kinds of event in one of several phrasings, and sometimes mentions a sector
word; each carries a reaction return equal to the event's planted effect plus noise of 3\% standard deviation. The first
60\% of the days (11\,923 headlines) train; the last 40\% (8\,077) test.

A headline such as ``Apex beat estimates but cut guidance amid copper demand'' passes through three steps before any model
sees it. The company mention is recognised and replaced by a placeholder, so that no model learns a company's name as a
signal; the text is lowercased and split into the tokens \texttt{<co> beat estimates but cut guidance amid copper
demand}; and the bigrams add \texttt{<co> beat}, \texttt{beat estimates}, \texttt{estimates but}, \texttt{but cut},
\texttt{cut guidance} and three more. The training headlines' 123 words (those found in at least three headlines) become 473
features with bigrams; in real news the vocabulary runs to tens of thousands and bigrams multiply it again.

\section{Sentiment from dictionaries and from returns}

A dictionary (Book~7, chapter~12) is a fixed instrument: the chapter's hand-built lists hold 11 positive and 16 negative
words. It reads a result correctly (``beat'' scores $+0.20$ on average, against a planted effect of $+1\%$; a guidance cut
$-0.18$ for $-1.5\%$) and misreads everything built from more than one word: ``beat estimates but cut guidance'' scores
exactly zero, one hit each way, where the planted effect is $-1.2\%$; ``did not cut guidance'' scores negative for a
planted $+0.3\%$; ``wins industry award'' scores positive for nothing; ``announces share buyback'' scores nothing for
$+0.7\%$. Loughran and McDonald (2011) showed that general word lists misclassify financial words; any list misses the
words its authors did not think of and cannot weigh the ones it has.

The alternative is to let returns choose the words. The return-supervised screen of \cref{lst:ml:text:supervised} keeps the
words whose headlines' mean return differs from the others' by more than three standard errors and weighs each by that
difference, in the spirit of Ke, Kelly and Xiu (2019), whose screening-and-topic model builds a sentiment score adapted
to return prediction. On the training half it keeps 41 words, from ``trims'' ($-1.43\%$) and ``halves'' to ``takeover''
($+1.79\%$); it also keeps ``to'' and ``be'', which carry information only because of how the corpus's takeover
headlines are phrased, the kind of word a stop list would have removed and a return screen rightly keeps.
\Cref{tab:ml:text:ic} compares the scores on the test half; \cref{fig:ml:text:learning} shows how each depends on the
number of labelled headlines.

\begin{table}[htbp]
\centering\small
\begin{tabular}{@{}lc@{}}
\toprule
score & rank IC with the reaction return\\
\midrule
dictionary (11 positive, 16 negative words) & 0.224\\
tf-idf words, \omterm{def:ml:linear-and-regularised-baselines:logistic}{logistic regression} & 0.273\\
tf-idf words and bigrams, \omterm{def:ml:linear-and-regularised-baselines:logistic}{logistic regression} & 0.257\\
return-supervised word list & 0.238\\
\omterm{def:ml:text-from-bag-of-words-to-embeddings:embedding}{word embeddings} (co-occurrence SVD), ridge regression & 0.094\\
event rules, mean training return per event & 0.292\\
\midrule
the planted effect & 0.295\\
\bottomrule
\end{tabular}
\caption{Rank information coefficients on the 8\,077 test headlines, models trained on the 11\,923 earlier ones. Data:
\texttt{ml\_text.scores}.}\label{tab:ml:text:ic}
\end{table}

\begin{omfigure}
\begin{tikzpicture}
\begin{semilogxaxis}[width=11cm,height=6.2cm,xlabel={labelled training headlines},ylabel={rank IC on the test half},
  tick label style={font=\small},label style={font=\small},xmin=250,xmax=12000,ymin=-0.05,ymax=0.32,
  xtick={300,1000,3000,10000},xticklabels={300,1\,000,3\,000,10\,000},grid=major,grid style={gray!25},legend columns=3,
  legend style={font=\footnotesize,at={(0.5,-0.25)},anchor=north,draw=none,fill=none,
  /tikz/every even column/.append style={column sep=6pt}},table/col sep=comma]
\addplot[black!60,thick,dashed] table[x=docs,y=truth]{figdata/ml/13-text-from-bag-of-words-to-embeddings/learning.csv};
\addplot[omThm,very thick,mark=*] table[x=docs,y=dictionary]{figdata/ml/13-text-from-bag-of-words-to-embeddings/learning.csv};
\addplot[omProp,very thick,mark=square*] table[x=docs,y=tfidf1]{figdata/ml/13-text-from-bag-of-words-to-embeddings/learning.csv};
\addplot[omProp!60,thick,mark=square] table[x=docs,y=tfidf2]{figdata/ml/13-text-from-bag-of-words-to-embeddings/learning.csv};
\addplot[omDef,very thick,mark=triangle*] table[x=docs,y=supervised]{figdata/ml/13-text-from-bag-of-words-to-embeddings/learning.csv};
\addplot[gray,very thick,mark=diamond*] table[x=docs,y=embeddings]{figdata/ml/13-text-from-bag-of-words-to-embeddings/learning.csv};
\addplot[black,thick,mark=o] table[x=docs,y=events]{figdata/ml/13-text-from-bag-of-words-to-embeddings/learning.csv};
\legend{planted effect,dictionary,tf-idf words,tf-idf with bigrams,supervised words,embeddings,event rules}
\end{semilogxaxis}
\end{tikzpicture}

\omcaption{Rank IC on the test half against the number of labelled headlines each model was trained on (the dictionary
uses none). With 300 headlines no word passes the supervised screen and its score is zero. Data:
\texttt{ml\_text.learning\_curve}.}\label{fig:ml:text:learning}
\end{omfigure}

Three lessons come out of the curves. A dictionary costs no labels: the tf-idf word model passes it only between 1\,000
and 3\,000 labelled headlines, the supervised word list and the bigram model only beyond 3\,000, and for real text, with a
far larger vocabulary and weaker signals, enough labels means years of news. The \omterm{def:ml:linear-and-regularised-baselines:logistic}{logistic regression} on tf-idf words (chapter~4's model, with a ridge penalty) ends
highest of the bag-of-words scores, at 0.273 against the dictionary's 0.224. And bigrams, which can represent the compound
headlines a word model cannot, cost more in variance than they gained: 0.257 against 0.273, with nearly four times the
features and the compounds only 9\% of the headlines. A richer representation needs more data to pay for itself.

\section{Topics and events}

\begin{definition}[Topic model, latent Dirichlet allocation]\label{def:ml:text-from-bag-of-words-to-embeddings:lda}
A \emph{topic model}\index{topic model} describes each document as a mixture of a few topics, each topic a distribution over
words, all estimated from word counts without labels. \emph{Latent Dirichlet allocation}\index{latent Dirichlet allocation}
(LDA) is the Bayesian topic model in which each document's topic shares and each topic's word probabilities have Dirichlet
priors, and each word is drawn by first drawing its topic from the document's shares (Blei, Ng and Jordan, 2002).
\end{definition}

\begin{table}[htbp]
\centering\small
\begin{tabular}{@{}cl@{}}
\toprule
topic & six most probable words\\
\midrule
1 & analysts, fell, short, expectations, forecasts, exceeded\\
2 & outlook, guidance, but, expectations, forecasts, estimates\\
3 & annual, meeting, record, attendance, reports, holds\\
4 & full, forecast, year, reports, line, costs\\
5 & wins, authorises, program, repurchase, industry, award\\
6 & announces, gets, rating, earnings, schedules, call\\
\bottomrule
\end{tabular}
\caption{Six topics fitted by LDA to the training headlines (the company placeholder, prepositions and context words
removed). Data: \texttt{ml\_text.topics}.}\label{tab:ml:text:topics}
\end{table}

\Cref{tab:ml:text:topics} shows what a \omterm{def:ml:text-from-bag-of-words-to-embeddings:lda}{topic model} finds in six-word headlines: the corpus's templates, not its news.
Beats and misses share topic~1, guidance changes and compounds of either direction share topic~2, and a buyback shares
topic~5 with an industry award; on
average only 63\% of an event's headlines fall in its most common topic. A topic is a pattern of co-occurrence, and words
of opposite meaning co-occur with the same words. Topics are useful to organise long documents (which sections of a
filing changed, what an earnings call spent its time on) and as features for a supervised model, not as a sentiment.

\begin{definition}[Event extraction]\label{def:ml:text-from-bag-of-words-to-embeddings:events}
\emph{Event extraction}\index{event extraction} turns text into structured records (who, what kind of event, which
direction, when), by rules on keywords and syntax or by a trained tagger, so that a model can be fitted on events rather
than on words.
\end{definition}

The chapter's extractor is a dozen rules: a keyword per event, ``not'' before a guidance cut, a result and a
guidance change joined by ``but''. It reads every one of the 20\,000 headlines correctly, and a score that assigns each
extracted event its mean training return reaches 0.292 against the planted effect's 0.295, above every bag-of-words model.
It should: its rules were written from the corpus's own templates. On a headline it was not written for (``surpasses
consensus'') it returns nothing. Rules are precise and brittle; learned models are robust and hungry for labels; production
systems use rules for the events that matter most and learned taggers for the rest, and measure both on text they were not
built from.

\section{Embeddings}

\begin{definition}[Word embedding]\label{def:ml:text-from-bag-of-words-to-embeddings:embedding}
A \emph{word embedding}\index{word embedding} maps each word to a dense vector of a few dozen to a few thousand real numbers,
learned so that words used in similar contexts get nearby vectors: by factorising a matrix of co-occurrence statistics, or
by training a network to predict a word from its neighbours (Mikolov and co-authors, 2013).
\end{definition}

The chapter's embeddings are the first kind (\cref{lst:ml:text:ppmi}): the positive pointwise mutual information of words
appearing in the same headline, reduced to 16 dimensions by a singular value decomposition (Book~4, chapter~25), rows
scaled to unit length. They learn exactly what the definition promises. Synonyms that fill the same slot, ``beat'' and
``tops'', ``raised'' and ``lifts'', ``guidance'' and ``outlook'', have cosine similarities above 0.99; so do antonyms that
fill the same slot, ``beat'' and ``missed'' at 0.95 (\cref{fig:ml:text:pairs}). Context decides similarity, and a word and
its opposite share their contexts. A ridge regression on the average of a headline's word vectors then has little to work
with: IC 0.094. Embeddings trained on a large general corpus carry more than this toy's, but the property is the same, and
it is why sentiment models built on embeddings are fine-tuned on labels (chapter~10) and why chapter~14's language models
are evaluated on the task, not on their vectors.

\begin{omfigure}
\begin{tikzpicture}
\begin{axis}[width=10cm,height=7cm,xbar,bar width=7pt,bar shift=0pt,area legend,xmin=0,xmax=1.05,
  xlabel={cosine similarity of the two word vectors},ytick={0,1,2,3,4,5,6,7,8,9,10},yticklabels from table={figdata/ml/13-text-from-bag-of-words-to-embeddings/pairs.csv}{pair},y dir=reverse,
  tick label style={font=\small},label style={font=\small},grid=major,grid style={gray!25},
  legend style={font=\footnotesize,at={(0.5,-0.18)},anchor=north,draw=none,fill=none},legend columns=4,
  table/col sep=comma,enlarge y limits=0.06]
\addplot[fill=omDef,draw=none,restrict expr to domain={\thisrow{i}}{0:2}] table[x=cosine,y=i]{figdata/ml/13-text-from-bag-of-words-to-embeddings/pairs.csv};
\addplot[fill=omProp,draw=none,restrict expr to domain={\thisrow{i}}{3:4}] table[x=cosine,y=i]{figdata/ml/13-text-from-bag-of-words-to-embeddings/pairs.csv};
\addplot[fill=omThm,draw=none,restrict expr to domain={\thisrow{i}}{5:8}] table[x=cosine,y=i]{figdata/ml/13-text-from-bag-of-words-to-embeddings/pairs.csv};
\addplot[fill=gray,draw=none,restrict expr to domain={\thisrow{i}}{9:10}] table[x=cosine,y=i]{figdata/ml/13-text-from-bag-of-words-to-embeddings/pairs.csv};
\legend{same meaning,same sector,opposite,unrelated}
\end{axis}
\end{tikzpicture}

\omcaption{Cosine similarities of \omterm{def:ml:text-from-bag-of-words-to-embeddings:embedding}{word embeddings} learned from the training headlines' co-occurrences. Opposites that fill
the same slot are as close as synonyms. Data: \texttt{ml\_text.pairs}.}\label{fig:ml:text:pairs}
\end{omfigure}

\section{Entity linking to securities}

\begin{definition}[Named-entity recognition, entity linking]\label{def:ml:text-from-bag-of-words-to-embeddings:ner}
\emph{Named-entity recognition}\index{named-entity recognition} finds the spans of a text that name entities (companies,
people, places, products) and their types. \emph{Entity linking}\index{entity linking} maps each recognised span to one
identifier in a knowledge base; for a trading firm, to a permanent identifier in the security master (Book~7, chapter~4) as
of the document's date.
\end{definition}

A model that reads news perfectly and trades the wrong company is worse than no model. The chapter's recogniser is a
gazetteer: every name, short name and ticker any company has held, matched longest first, which finds every mention in the
corpus because the corpus uses no other names; real recognisers also tag names never seen before, from capitalisation,
syntax and a trained tagger. The linker (\cref{lst:ml:text:link}) then resolves a ticker through the security master and a
name through an alias table, both as of the headline's date, and separates several candidates by sector words in the
headline, then by coverage.

\begin{table}[htbp]
\centering\small
\begin{tabular}{@{}lcccccc@{}}
\toprule
& \multicolumn{3}{c}{aliases as of the headline's date} & \multicolumn{3}{c}{today's aliases}\\
\cmidrule(lr){2-4}\cmidrule(l){5-7}
mention & correct & linked & precision & correct & linked & precision\\
\midrule
full name & 100.0 & 100.0 & 100.0 & 89.3 & 89.3 & 100.0\\
short name & 98.2 & 100.0 & 98.2 & 86.8 & 88.6 & 98.0\\
ticker & 100.0 & 100.0 & 100.0 & 84.5 & 92.4 & 91.5\\
\midrule
all & 99.5 & 100.0 & 99.5 & 87.5 & 89.7 & 97.5\\
\bottomrule
\end{tabular}
\caption{\omterm{def:ml:text-from-bag-of-words-to-embeddings:ner}{Entity linking} on the 20\,000 headlines (\%): share linked to the right company, share linked to any company, and
precision (right among linked). Data: \texttt{ml\_text.linking}.}\label{tab:ml:text:linking}
\end{table}

\Cref{tab:ml:text:linking} holds the chapter's second result. With point-in-time aliases the linker is right 99.5\% of the
time; its only errors are the short name Apex in the 245 headlines that carry no sector word, where it falls back on the
most covered Apex and is right 58.0\% of the time (with a sector word it is right in all 737). With today's aliases it is
right 87.5\% of the time: old names and old tickers resolve to nothing, and the four reused tickers resolve to the new
listing, so that 8.5\% of ticker mentions it links go to the wrong company, silently. A backtest on those links would trade
a takeover headline about one company in the shares of another; the rule of Book~7, chapter~3, applies to aliases as to
prices.

\begin{method}[A text pipeline for trading]\label{met:ml:text:method}
\begin{enumerate}
\item Recognise mentions and link them point in time to permanent identifiers; measure linking precision on a labelled
sample, by mention kind.
\item Mask the mentions, tokenise, and build the representation (tf-idf words first).
\item Start with a dictionary as the baseline; fit supervised scores on returns in walk-forward folds, with the reaction
window and horizon fixed in advance.
\item Extract the events that matter with rules or a tagger and test them against the bag-of-words model.
\item Treat embeddings and topics as features to be validated, not as sentiment.
\end{enumerate}
\end{method}

\section{Tutorial: which Apex?}\label{tut:ml:text-from-bag-of-words-to-embeddings:tut}

\begin{tutorial}
\textbf{Goal.} Generate the corpus, recognise and link mentions point in time, and compare five text scores by their
information coefficient. \textbf{End state:} \cref{tab:ml:text:ic,tab:ml:text:topics,tab:ml:text:linking},
\cref{fig:ml:text:learning,fig:ml:text:pairs}.
\begin{tutsteps}
\item \textbf{A return-supervised word list.}
\omcode{code/firm/textml/firm_textml.py}{51}{70}{A return-supervised word screen.}\label{lst:ml:text:supervised}
\item \textbf{Embeddings from co-occurrence.}
\omcode{code/firm/textml/firm_textml.py}{78}{99}{Word embeddings from positive pointwise mutual information and an SVD.}\label{lst:ml:text:ppmi}
\item \textbf{Point-in-time linking.}
\omcode{code/firm/textml/firm_textml.py}{211}{222}{Linking a mention to a permanent identifier as of a date.}\label{lst:ml:text:link}
\item \textbf{Run} \texttt{ml\_text.scores()}, \texttt{learning\_curve()}, \texttt{topics()}, \texttt{pairs()},
\texttt{linking()} and \texttt{fig\_text.py}.
\end{tutsteps}
\textbf{What to change next.} Add phrasings to the test period that the training period never saw, and watch the rules and
the supervised models fail differently; give the linker a context model instead of a sector-word list.
\end{tutorial}

\section{Build: text features and entity linking}\label{bld:ml:text-from-bag-of-words-to-embeddings:textml}

\begin{build}
\bfield{Purpose} Text turned into scores that can be validated like any other feature, attached to the right security.
\bfield{Interface} \texttt{tokenize}, \texttt{ngrams}, \texttt{dictionary\_score(tokens, pos, neg)},
\texttt{supervised\_words(docs, returns, min\_count, z)}, \texttt{supervised\_score}, \texttt{ppmi\_svd(docs, k,
min\_count)}, \texttt{doc\_vectors}, \texttt{extract\_event(tokens)}; \texttt{AliasTable}, \texttt{recognise(text,
gazetteer)}, \texttt{link(surface, date, context, table, sm, sectors, prior)}; \texttt{news\_corpus(n\_docs, seed, days)}
on \texttt{firm.secmaster}.
\bfield{Rules} Mentions are masked before any model sees the text; aliases and tickers are resolved as of the document's
date; supervised scores are fitted on earlier documents than they score.
\bfield{Acceptance tests} \texttt{code/firm/textml/tests/}: tokens and \omterm{def:ml:text-from-bag-of-words-to-embeddings:bow}{n-grams} of a known headline; the supervised screen
finds a planted word and ignores a neutral one; embeddings put words with identical contexts together; the rules read
compounds and negations; point-in-time linking follows a ticker change, a rename and a reused ticker, and today's table
does not.
\bfield{Stretch} A trained sequence tagger for recognition; linking by a context model over the candidates' descriptions;
the Ke--Kelly--Xiu \omterm{def:ml:text-from-bag-of-words-to-embeddings:lda}{topic model} on top of the screen.
\end{build}

\begin{omsources}
\item P.~C.~Tetlock, ``Giving content to investor sentiment: the role of media in the stock market'', \emph{Journal of
Finance} 62(3), 2007.
\item T.~Loughran and B.~McDonald, ``When is a liability not a liability? Textual analysis, dictionaries, and 10-Ks'',
\emph{Journal of Finance} 66(1), 2011.
\item Z.~T.~Ke, B.~T.~Kelly and D.~Xiu, ``Predicting returns with text data'', NBER Working Paper 26186, 2019.
\item M.~Gentzkow, B.~Kelly and M.~Taddy, ``Text as data'', \emph{Journal of Economic Literature} 57(3), 2019.
\item D.~M.~Blei, A.~Y.~Ng and M.~I.~Jordan, ``Latent Dirichlet allocation'', \emph{Advances in Neural Information Processing
Systems} 14, MIT Press, 2002.
\item T.~Mikolov, K.~Chen, G.~Corrado and J.~Dean, ``Efficient estimation of word representations in vector space'',
arXiv:1301.3781, 2013.
\end{omsources}

\section{Exercises}

\begin{exercise}[$\star$]\label{exo:ml:text-from-bag-of-words-to-embeddings:1}
List the tokens and the bigrams of ``\texttt{<co>} did not cut guidance''. Which bigram carries the negation?
\end{exercise}

\begin{exercise}[$\star$]\label{exo:ml:text-from-bag-of-words-to-embeddings:2}
In a corpus of 20\,000 documents, ``guidance'' appears in 3\,000 and ``halves'' in 200. What are their inverse document
frequencies ($\log N/n_w$)? Which weighs more in a document containing both once?
\end{exercise}

\begin{exercise}[$\star$]\label{exo:ml:text-from-bag-of-words-to-embeddings:3}
Compute the dictionary score of ``\texttt{<co>} beat estimates but cut guidance'' (six tokens) and of ``\texttt{<co>} wins
industry award'' (four tokens) with the chapter's lists. What are the planted effects?
\end{exercise}

\begin{exercise}[$\star\star$]\label{exo:ml:text-from-bag-of-words-to-embeddings:4}
Why do ``beat'' and ``missed'' get almost the same embedding here, and what would separate them?
\end{exercise}

\begin{exercise}[$\star\star$]\label{exo:ml:text-from-bag-of-words-to-embeddings:5}
The supervised screen keeps ``to'' and ``be''. Is that a bug? When would it become one?
\end{exercise}

\begin{exercise}[$\star\star$]\label{exo:ml:text-from-bag-of-words-to-embeddings:6}
\emph{Find the flaw.} ``We linked five years of news to tickers with the vendor's current symbol file; 97\% of the links we
checked by hand were right, so linking is not a problem.''
\end{exercise}

\begin{exercise}[$\star\star\star$]\label{exo:ml:text-from-bag-of-words-to-embeddings:7}
\emph{Coding.} Train the supervised screen and the tf-idf \omterm{def:ml:linear-and-regularised-baselines:logistic}{logistic regression} on the training half with the company
mentions left in the text (unmasked). What do the models learn about the companies, and what happens to the test IC? Explain.
\end{exercise}

\begin{exercise}[$\star\star\star$]\label{exo:ml:text-from-bag-of-words-to-embeddings:8}
Show that with a document's tokens as a \omterm{def:ml:text-from-bag-of-words-to-embeddings:bow}{bag of words}, no linear model on word counts can give ``beat \dots\ but cut'' an
effect of $-1.2$ and ``missed \dots\ but raised'' an effect of $+1.2$ while giving ``beat'', ``missed'', ``raised'' and ``cut''
alone $+1$, $-1$, $+1.5$ and $-1.5$. What features would?
\end{exercise}

\section{Problem: Which Apex?}

\begin{problem}[{Weekend problem --- text to the right trade}]\label{pb:ml:text-from-bag-of-words-to-embeddings:1}
The chapter's corpus: 20\,000 headlines, 32 companies (three called Apex), six ticker changes, three renames, four reused
tickers.

\medskip\noindent\textbf{Part I --- Representation.}
\begin{enumerate}
\item Why mask the company mentions before building features?
\item What does tf-idf do to a word that appears in every headline?
\item How many features do bigrams add, roughly, and why did they not help here?
\item What does the planted effect's IC of 0.295 say about the best any text model can do?
\end{enumerate}

\medskip\noindent\textbf{Part II --- Scores.}
\begin{enumerate}[resume]
\item Rank the six scores of \cref{tab:ml:text:ic}.
\item Below how many labelled headlines is the dictionary the best score?
\item Why do the event rules beat every learned model, and why should you not believe it?
\item Why do embeddings score 0.094?
\end{enumerate}

\medskip\noindent\textbf{Part III --- Linking.}
\begin{enumerate}[resume]
\item What share of mentions does point-in-time linking get right, and where are its errors?
\item What does today's table get right, and what are its two failure modes?
\item Which failure is worse for a backtest, an unlinked mention or a wrongly linked one?
\item How accurate is the linker on Apex with and without a sector word?
\end{enumerate}

\medskip\noindent\textbf{Part IV --- The verdict.}
\begin{enumerate}[resume]
\item State the \emph{named result}: entity-linking accuracy and precision with and without point-in-time aliases, and the
IC of the return-supervised models against the dictionary's.
\item What would you build first for a news desk: better sentiment or better linking?
\item How would you check a vendor's entity links?
\item How do you keep a supervised text model honest in time?
\item What is a \omterm{def:ml:text-from-bag-of-words-to-embeddings:lda}{topic model} good for in this business?
\item What changes with long documents (filings, transcripts) rather than headlines?
\item What would make the dictionary competitive again?
\item In one sentence: what is the most important step between a headline and a trade?
\end{enumerate}
\end{problem}

\section{Interview questions}

\begin{interviewq}[$\star$ \iqroles{researcher, mle}]\label{iq:ml:text-from-bag-of-words-to-embeddings:1}
What is tf-idf, and why does it help a linear model on text?
\end{interviewq}

\begin{interviewq}[$\star$ \iqroles{researcher}]\label{iq:ml:text-from-bag-of-words-to-embeddings:2}
Why do finance-specific dictionaries exist?
\end{interviewq}

\begin{interviewq}[$\star\star$ \iqroles{researcher}]\label{iq:ml:text-from-bag-of-words-to-embeddings:3}
How would you build a sentiment score supervised by returns, and what could leak into it?
\end{interviewq}

\begin{interviewq}[$\star\star$ \iqroles{mle}]\label{iq:ml:text-from-bag-of-words-to-embeddings:4}
Why are antonyms close in word-embedding space, and what does it mean for using embeddings as features?
\end{interviewq}

\begin{interviewq}[$\star\star$ \iqroles{mle, dev}]\label{iq:ml:text-from-bag-of-words-to-embeddings:5}
Design the entity-linking step for a news feed that drives trading. What can go wrong?
\end{interviewq}

\begin{interviewq}[$\star\star\star$ \iqroles{researcher}]\label{iq:ml:text-from-bag-of-words-to-embeddings:6}
Your text model's IC fell by half in its first live year. List the causes you would check, in order.
\end{interviewq}
